\documentclass[12pt]{article}
\usepackage{pst-plot}
\usepackage{fontspec}
\usepackage{graphicx}
\usepackage{scrextend}
\usepackage{listings}
\usepackage{tikz}
\usepackage{amssymb}
\usepackage{color}
\usepackage{soul}
\usepackage{fancyhdr}
\usepackage{pgfornament}
\usepackage[many]{tcolorbox}
\usepackage{collectbox}
\usepackage{mdframed}
\usepackage{xcolor}
\usepackage{mathtools}
\usepackage[hidelinks]{hyperref}
\usepackage[final]{pdfpages}
\usepackage[left=0.5in,right=0.5in,top=1in,bottom=1in,footskip=15mm]{geometry}

% ------------
% Custom fonts
% ------------
\newfontfamily{\quicksand}[Path=./fonts/,]{Quicksand-Regular}
\newfontfamily{\cousine}[Path=./fonts/,]{Cousine-Regular}
\newfontfamily{\cousinebd}[Path=./fonts/,]{Cousine-Bold}

% -------------
% Custom colors
% -------------
\definecolor{eBrightWhite}{RGB}{200,200,200}
\definecolor{eBlueOne}{RGB}{2,0,50}
\definecolor{ePurple}{RGB}{157,27,249}
\definecolor{eGreen}{RGB}{0,255,128}

% ---------------
% Custom commands
% ---------------
\newcommand{\sectionHeader}[3]{
\noindent {\Huge {\quicksand \textcolor{ePurple}{#1:} \hspace{3mm}\textcolor{white}{#2}}}

\vspace{15mm}
}

\linespread{1.2}

\renewcommand{\headrulewidth}{0pt}

\begin{document}
\pagecolor{eBlueOne}
\sectionHeader{Section 1}{Number Theory}

\begin{addmargin}[0em]{2em}
{\small \textcolor{eBrightWhite}{{\cousine {\cousinebd 1.} Circle the numbers in the table below that are divisible by 3.}}}

\vspace{10mm}

\begin{center}
{\cousine
\begin{tabular}{ccc}
\textcolor{white}{699630} \hspace{5mm} & \textcolor{white}{6661} \hspace{5mm} & \textcolor{white}{5052} \\
\\
\textcolor{white}{777} \hspace{5mm} & \textcolor{white}{9991332} \hspace{5mm} & \textcolor{white}{40000044} \\
\\
\textcolor{white}{66666} \hspace{5mm} & \textcolor{white}{783209} \hspace{5mm} & \textcolor{white}{0} \\
\end{tabular}
}
\end{center}

\end{addmargin}

\vspace{15mm}

\begin{addmargin}[0em]{2em}
{\small \textcolor{eBrightWhite}{{\cousine {\cousinebd 2.} Circle the numbers in the table below that are divisible by 4.}}}

\vspace{10mm}

\begin{center}
{\cousine
\begin{tabular}{ccc}
\textcolor{white}{7840} \hspace{5mm} & \textcolor{white}{1000} \hspace{5mm} & \textcolor{white}{5748324} \\
\\
\textcolor{white}{6661} \hspace{5mm} & \textcolor{white}{5423} \hspace{5mm} & \textcolor{white}{80000052} \\
\\
\textcolor{white}{5432} \hspace{5mm} & \textcolor{white}{4325} \hspace{5mm} & \textcolor{white}{0} \\
\end{tabular}
}
\end{center}

\end{addmargin}

\vspace{15mm}

\begin{addmargin}[0em]{2em}
{\small \textcolor{eBrightWhite}{{\cousine {\cousinebd 3.} Write a 10-digit number that is divisible by 12. The digits cannot all be the same number.}}}
\end{addmargin}

\vspace{15mm}

\begin{addmargin}[0em]{2em}
{\small \textcolor{eBrightWhite}{{\cousine {\cousinebd 4.} Why is it impossible to create a number that is divisible by 4 with only odd digits?}}}
\end{addmargin}

\vspace{15mm}

\begin{addmargin}[0em]{2em}
{\small \textcolor{eBrightWhite}{{\cousine {\cousinebd 5.} Order the digits 1, 9, and 2 to create a number that is divisible by 3, but not by 12.}}}
\end{addmargin}

\pagebreak
\newpage

\sectionHeader{Section 2}{Real Algebra}

\begin{addmargin}[0em]{2em}
{\small \textcolor{eBrightWhite}{{\cousine {\cousinebd 1.} Classify and solve {\Large $$4x + 5 = 21.$$}}}}
\end{addmargin}

\vspace{60mm}

\begin{addmargin}[0em]{2em}
{\small \textcolor{eBrightWhite}{{\cousine {\cousinebd 2.} Classify and solve {\Large $$\frac{1}{x-2}+7=10.$$}}}}
\end{addmargin}

\vspace{70mm}

\begin{addmargin}[0em]{2em}
{\small \textcolor{eBrightWhite}{{\cousine {\cousinebd 3.} Classify and solve {\Large $$x^2-7x-78=0.$$}}}}
\end{addmargin}

\sectionHeader{Section 3}{Complex Arithmetic}

\begin{addmargin}[0em]{2em}
{\small \textcolor{eBrightWhite}{{\cousine {\cousinebd 1.} {\Large $$7\cdot(7+6i)=$$}}}}
\end{addmargin}
\begin{addmargin}[0em]{2em}
{\small \textcolor{eBrightWhite}{{\cousine {\cousinebd 2.} {\Large $$(2+6i)^{*}=$$}}}}
\end{addmargin}
\begin{addmargin}[0em]{2em}
{\small \textcolor{eBrightWhite}{{\cousine {\cousinebd 3.} {\Large $$(2-8i)\cdot(8+7i)=$$}}}}
\end{addmargin}

\vspace{30mm}

\begin{addmargin}[0em]{2em}
{\small \textcolor{eBrightWhite}{{\cousine {\cousinebd 4.} {\Large $$(9+2i)\cdot(9-2i)=$$}}}}
\end{addmargin}
\begin{addmargin}[0em]{2em}
{\small \textcolor{eBrightWhite}{{\cousine {\cousinebd 5.} {\Large $$\frac{2+3i}{1-2i}=$$}}}}
\end{addmargin}

\pagebreak
\newpage

\sectionHeader{Section 4}{Coordinate Geometry}

\begin{addmargin}[0em]{2em}
{\small \textcolor{eBrightWhite}{{\cousine {\cousinebd 1.} Write the function that each graph corresponds to.}}}

\begin{tabular}{ccc}

\begin{pspicture}(-1.5,0)(0,3.5)
    \def\f(#1){#1^2}
    \psaxes[linecolor=white,tickcolor=white,labels=none,ticksize=-2pt 2pt]{->}(0,0)(-2,-2)(2,2)[\textcolor{white}{$x$},0][\textcolor{white}{$y$},90]
    \psplot[linecolor=eGreen,algebraic]{-1.5}{1.5}{\f(x)}
\end{pspicture} \hspace{40mm}

&

\begin{pspicture}(-1.5,0)(0,3.5)
    \def\f(#1){#1}
    \psaxes[linecolor=white,tickcolor=white,labels=none,ticksize=-2pt 2pt]{->}(0,0)(-2,-2)(2,2)[\textcolor{white}{$x$},0][\textcolor{white}{$y$},90]
    \psplot[linecolor=eGreen,algebraic]{-1.5}{1.5}{\f(x)}
\end{pspicture} \hspace{40mm}

&

\begin{pspicture}(-1.5,0)(0,3.5)
    \def\f(#1){#1^3}
    \psaxes[linecolor=white,tickcolor=white,labels=none,ticksize=-2pt 2pt]{->}(0,0)(-2,-2)(2,2)[\textcolor{white}{$x$},0][\textcolor{white}{$y$},90]
    \psplot[linecolor=eGreen,algebraic]{-1.25}{1.25}{\f(x)}
\end{pspicture} \\
\\
\\
\\
\\
\\
\begin{pspicture}(-1.5,0)(0,3.5)
    \def\f(#1){sin(#1)}
    \psaxes[linecolor=white,tickcolor=white,labels=none,ticksize=-2pt 2pt]{->}(0,0)(-2,-2)(2,2)[\textcolor{white}{$x$},0][\textcolor{white}{$y$},90]
    \psplot[linecolor=eGreen,algebraic]{-2.5}{2.5}{\f(x)}
\end{pspicture} \hspace{40mm}

&

\begin{pspicture}(-1.5,0)(0,3.5)
    \def\f(#1){abs(#1)}
    \psaxes[linecolor=white,tickcolor=white,labels=none,ticksize=-2pt 2pt]{->}(0,0)(-2,-2)(2,2)[\textcolor{white}{$x$},0][\textcolor{white}{$y$},90]
    \psplot[linecolor=eGreen,algebraic]{-2}{2}{\f(x)}
\end{pspicture} \hspace{40mm}

&

\begin{pspicture}(-1.5,0)(0,3.5)
    \def\f(#1){log(#1)}
    \psaxes[linecolor=white,tickcolor=white,labels=none,ticksize=-2pt 2pt]{->}(0,0)(-2,-2)(2,2)[\textcolor{white}{$x$},0][\textcolor{white}{$y$},90]
    \psplot[linecolor=eGreen,algebraic]{0.01}{3}{\f(x)}
\end{pspicture} \\
\\
\\
\\
\\
\\
\begin{pspicture}(-1.5,0)(0,3.5)
    \def\f(#1){cos(#1)}
    \psaxes[linecolor=white,tickcolor=white,labels=none,ticksize=-2pt 2pt]{->}(0,0)(-2,-2)(2,2)[\textcolor{white}{$x$},0][\textcolor{white}{$y$},90]
    \psplot[linecolor=eGreen,algebraic]{-2.5}{2.5}{\f(x)}
\end{pspicture} \hspace{40mm}

&

\begin{pspicture}(-1.5,0)(0,3.5)
    \def\f(#1){1/#1}
    \psaxes[linecolor=white,tickcolor=white,labels=none,ticksize=-2pt 2pt]{->}(0,0)(-2,-2)(2,2)[\textcolor{white}{$x$},0][\textcolor{white}{$y$},90]
    \psplot[linecolor=eGreen,algebraic]{0.5}{2}{\f(x)}
    \psplot[linecolor=eGreen,algebraic]{-2}{-0.5}{\f(x)}
\end{pspicture} \hspace{40mm}

&

\begin{pspicture}(-1.5,0)(0,3.5)
    \def\f(#1){2^#1}
    \psaxes[linecolor=white,tickcolor=white,labels=none,ticksize=-2pt 2pt]{->}(0,0)(-2,-2)(2,2)[\textcolor{white}{$x$},0][\textcolor{white}{$y$},90]
    \psplot[linecolor=eGreen,algebraic]{-2}{1.5}{\f(x)}
\end{pspicture} \\
\\
\\
\\
\\

\end{tabular}

\end{addmargin}

\pagebreak
\newpage

\sectionHeader{Section 5}{Calculus}

\begin{addmargin}[0em]{2em}
{\small \textcolor{eBrightWhite}{{\cousine {\cousinebd 1.} Find the derivative of {\Large $$f(x) = \sin(\tan(\sec(\cos(x)))).$$}}}}
\end{addmargin}

\vspace{40mm}

\begin{addmargin}[0em]{2em}
{\small \textcolor{eBrightWhite}{{\cousine {\cousinebd 2.} Find the derivative of {\Large $$f(x) = x^{2022}.$$}}}}
\end{addmargin}

\vspace{40mm}

\begin{addmargin}[0em]{2em}
{\small \textcolor{eBrightWhite}{{\cousine {\cousinebd 3.} Solve {\Large $$\int \sec^2(x) \hspace{2mm} dx.$$}}}}
\end{addmargin}

\pagebreak
\newpage

\sectionHeader{Section 6}{Challenge}

\begin{addmargin}[0em]{2em}
{\small \textcolor{eBrightWhite}{{\cousine {\cousinebd 1.} Solve {\Large $$\int 2xe^{x^2} \hspace{2mm} dx.$$}}}}
\end{addmargin}

\end{document}
